\chapter{KESIMPULAN DAN SARAN}

% Ubah konten-konten berikut sesuai dengan yang ingin diisi pada bab ini

\section{Kesimpulan}

Berdasarkan seluruh rangkaian kegiatan kerja praktik yang telah dilakukan mulai dari perancangan hingga pengujian sistem \textit{fault tolerance} untuk mitigasi artefak gerak pada citra MRI otak, dapat ditarik beberapa kesimpulan utama sebagai berikut.

\begin{enumerate}[nolistsep]

  \item Simulasi artefak gerak berupa translasi dan rotasi berhasil dilakukan dengan memanipulasi data citra pada domain frekuensi (\textit{k-space}) untuk menghasilkan dataset pelatihan sintetis yang realistis dan meniru inkonsistensi pada mesin MRI.

  \item Pengembangan model klasifikasi menggunakan strategi transfer learning menunjukkan bahwa arsitektur berbasis Transformer, terutama SwinV2 Transformer, memiliki keunggulan dibandingkan arsitektur CNN dalam menangkap pola artefak yang kompleks baik secara lokal maupun global.

  \item Modul koreksi artefak berhasil dikembangkan menggunakan arsitektur Pix2Pix GAN yang mampu memulihkan detail anatomi sekaligus menghilangkan pola \textit{ghosting} dan \textit{blurring} pada citra dengan memanfaatkan fungsi \textit{loss} gabungan.
  
  \item Hasil evaluasi menunjukkan kinerja sistem yang sangat baik dimana modul klasifikasi terbaik mencapai akurasi 93,82\% dan hasil generasi gambar modul koreksi menghasilkan kenaikan skor SSIM dibandingkan dengan gambar dengan artefak meskipun skor PSNR menunjukkan penurunan, yang menunjukkan bahwa model berfokus pada pemulihan fitur-fitur penting yang meningkatkan kualitas visual gambar dibandingkan dengan perbaikan kualitas pada level piksel.

\end{enumerate}

\section{Saran}

Untuk meningkatkan kinerja dan memperluas cakupan sistem ini di masa mendatang, terdapat beberapa langkah strategis yang dapat dikembangkan lebih lanjut.

\begin{enumerate}[nolistsep]

  \item Pengembangan selanjutnya disarankan dapat mencakup jenis artefak medis lainnya yang sering ditemui dalam MRI seperti artefak \textit{aliasing} atau \textit{noise} frekuensi tinggi agar sistem memiliki cakupan proteksi yang lebih menyeluruh.

  \item Transisi dari pemrosesan irisan dua dimensi menuju pemrosesan volume MRI tiga dimensi secara utuh sangat disarankan guna menjaga konsistensi fitur anatomi antar-irisan yang lebih baik.

  \item Validasi menggunakan dataset artefak dari mesin MRI di rumah sakit diperlukan untuk menguji kemampuan generalisasi model terhadap gerakan pasien yang sebenarnya di lingkungan klinis yang lebih dinamis.
  
  \item Pengujian integrasi penuh secara \textit{end-to-end} dengan sistem \textit{image captioning} medis perlu dilakukan untuk mengukur peningkatan kualitas deskripsi klinis yang dihasilkan secara langsung setelah melalui proses koreksi artefak.

\end{enumerate}
